\documentclass[a4paper,11pt,dvipdfmx]{jsarticle}

\usepackage[top=20mm, bottom=25mm, left=20mm, right=20mm]{geometry}
\usepackage{url}
\usepackage{graphicx}
\usepackage{amssymb}
\usepackage{amsmath}
\usepackage{booktabs}
\usepackage{float}
\usepackage{multirow}
\usepackage{tikz}
\usepackage{pgfplots}
\usetikzlibrary{positioning}
\pgfplotsset{compat=1.18}

\title{今月の進捗報告\\[0.5em]\large 前年の月内最大買電電力が年間電気料金に与える影響}
\author{神戸大学大学院システム情報学研究科 修士1年\\山崎博之}
\date{2026年7月}

\begin{document}

\maketitle

\section{背景と目的}

高圧電力契約の基本料金は，各月の契約電力に基本料金単価を乗じて算定される．契約電力算定では，当月を含む直近12か月それぞれの月内最大買電電力を比較し，そのうち最も大きい値を当月の契約電力とする．このため，2026年の蓄電池運用を決める際には，2026年の需要とPV発電量だけでなく，2025年各月の月内最大買電電力も基本料金に影響する．

たとえば，2025年12月の月内最大買電電力は，2026年1月から11月までの契約電力算定に含まれ，2026年12月の算定から外れる．したがって，同じ大きさの月内最大買電電力であっても，それが記録された月によって，2026年の契約電力算定に含まれ続ける期間が異なる．

今回は，2025年の月内最大買電電力の大きさと，その値が2026年の契約電力算定に含まれ続ける月数を変化させ，年間基本料金，年間電力量料金および年間電気料金に与える影響を調査した．これにより，過去の月内最大買電電力の大きさだけでなく，契約電力算定に残る期間を考慮する必要性を明らかにする．

\section{対象期間と集合}

運用計画の対象期間は2026年1月1日0時0分～12月31日23時30分とし，時間間隔は $\Delta t=0.5$\,hとする．需要電力とPV発電電力には，既存の2025年の30分値データを2026年の同月・同時刻に対応させて用いる．これらの値と電力量料金単価は，すべての30分時間帯に対して既知であり，1年間の蓄電池充放電計画を一括で決定する．

次の集合を定義する．

\begin{itemize}
    \item $\mathcal{M}=\{1,\ldots,12\}$：年によらない月番号の集合．最適化対象年である2026年の月の添字に $m$ を用いる．
    \item $\mathcal{K}=\{0,\ldots,17519\}$：2026年の30分時間帯の集合．時間帯の添字に $k$ を用いる．
    \item $\mathcal{K}_m \subset \mathcal{K}$：2026年の月 $m$ に属する30分時間帯の集合．
    \item $\mathcal{M}_m^{2025}=\{j\in\mathcal{M}\mid m<j\}$：2026年の月 $m$ の契約電力算定に含まれる，2025年の月番号の集合．
    \item $\mathcal{M}_m^{2026}=\{j\in\mathcal{M}\mid j\leq m\}$：2026年の月 $m$ の契約電力算定に含まれる，2026年の月番号の集合．
\end{itemize}
% 「SOCを表す時点の集合」ってどういう意味

たとえば，2026年1月の契約電力は2025年2月から12月と2026年1月を参照するため，$\mathcal{M}_1^{2025}=\{2,\ldots,12\}$，$\mathcal{M}_1^{2026}=\{1\}$ である．2026年12月には2025年の値を参照せず，$\mathcal{M}_{12}^{2025}=\emptyset$，$\mathcal{M}_{12}^{2026}=\mathcal{M}$ となる．

\section{運用手法}

\subsection{記号の定義}

本報告で用いる主な記号を表\ref{tab:symbols}に示す．対象システムと充放電電力の対応を図\ref{fig:system_flow}に示す．

\begin{figure}[H]
    \centering
    \begin{tikzpicture}[
            block/.style={rectangle,draw,minimum width=27mm,minimum height=9mm,align=center},
            flow/.style={->,thick}
        ]
        \node[circle,fill,inner sep=1.5pt] (bus) {};
        \node[block,left=30mm of bus] (grid) {系統電力};
        \node[block,above=18mm of bus] (pv) {PV};
        \node[block,right=30mm of bus] (load) {施設需要};
        \node[block,below=18mm of bus] (converter) {電力変換器};
        \node[block,below=11mm of converter] (battery) {蓄電池};

        \draw[flow] (grid) -- node[above] {$s_k^{\mathrm{BY}}$} (bus);
        \draw[flow] (pv) -- node[right] {$g_k^{\mathrm{P2}}$} (bus);
        \draw[flow] (bus) -- node[above] {$d_k$} (load);
        \draw[flow] ([xshift=-3mm]bus.south) -- node[left] {$x_k^{\mathrm{C1}}$} ([xshift=-3mm]converter.north);
        \draw[flow] ([xshift=3mm]converter.north) -- node[right] {$x_k^{\mathrm{D2}}$} ([xshift=3mm]bus.south);
        \draw[flow] ([xshift=-3mm]converter.south) -- node[left] {$x_k^{\mathrm{C2}}$} ([xshift=-3mm]battery.north);
        \draw[flow] ([xshift=3mm]battery.north) -- node[right] {$x_k^{\mathrm{D1}}$} ([xshift=3mm]converter.south);
    \end{tikzpicture}
    \caption{対象システムの電力フローと充放電電力の定義}
    \label{fig:system_flow}
\end{figure}

\begin{table}[H]
    \centering
    \small
    \caption{主な記号}
    \label{tab:symbols}
    \begin{tabular}{c p{92mm} c}
        \toprule
        記号                                            & 定義                            & 単位     \\
        \midrule
        \multicolumn{3}{l}{決定変数}                                                               \\
        $s_k^{\mathrm{BY}}$                           & 時間帯 $k$ の30分平均買電電力            & kW     \\
        $g_k^{\mathrm{P2}}$                           & 実際に使用するPV発電電力                 & kW     \\
        $x_k^{\mathrm{C1}},x_k^{\mathrm{C2}}$         & 充電電力（変換前／変換後）                 & kW     \\
        $x_k^{\mathrm{D1}},x_k^{\mathrm{D2}}$         & 放電電力（変換前／変換後）                 & kW     \\
        $b_k$                                         & 時点 $k$ の蓄電池SOC                & kWh    \\
        $z_k$                                         & 充放電モード（$1$：充電，$0$：放電）         & --     \\
        $s_m^{2026}$                                  & 2026年の月 $m$ における月内最大買電電力      & kW     \\
        $s_m^{\mathrm{CON}}$                          & 2026年の月 $m$ に適用する契約電力         & kW     \\
        \midrule
        \multicolumn{3}{l}{既知パラメータ}                                                            \\
        $d_k$                                         & 需要電力                          & kW     \\
        $g_k^{\mathrm{P1}}$                           & 利用可能なPV発電電力                   & kW     \\
        $s_m^{2025}$                                  & 2025年の月 $m$ における月内最大買電電力（実績値） & kW     \\
        $p^{\mathrm{BY}}$                             & 買電電力量単価                       & 円/kWh  \\
        \midrule
        \multicolumn{3}{l}{定数}                                                                 \\
        $\Delta t$                                    & 1ステップの長さ（$=0.5$）              & h      \\
        $\alpha$                                      & 充電および放電効率（$=0.98$）            & --     \\
        $C$                                           & 蓄電池定格容量（$=860$）               & kWh    \\
        $P_{\max}^{\mathrm{C}},P_{\max}^{\mathrm{D}}$ & 最大充電および放電出力（$=400$）           & kW     \\
        $p^{\mathrm{B}}$                              & 力率補正後の月額基本料金単価（$=2{,}405.16$） & 円/kW/月 \\
        \bottomrule
    \end{tabular}

\end{table}

\subsection{蓄電池運用の制約}

各 $k\in\mathcal{K}$ における電力需給バランスを

\begin{equation}
    g_k^{\mathrm{P2}}+s_k^{\mathrm{BY}}+x_k^{\mathrm{D2}}
    =d_k+x_k^{\mathrm{C1}},
    \quad \forall k\in\mathcal{K}
\end{equation}

とする．PV使用電力は

\begin{equation}
    0\leq g_k^{\mathrm{P2}}\leq g_k^{\mathrm{P1}},
    \quad \forall k\in\mathcal{K}
\end{equation}

を満たす．また，売電を行わないため，買電電力，月内最大買電電力および契約電力は非負とする．

\begin{equation}
    s_k^{\mathrm{BY}}\geq0,\quad
    s_m^{2026}\geq0,\quad s_m^{\mathrm{CON}}\geq0,
    \quad \forall k\in\mathcal{K},\ \forall m\in\mathcal{M}.
\end{equation}

充電および放電効率を $\alpha=0.98$ とし，変換前後の電力の関係を

\begin{align}
    x_k^{\mathrm{C2}} & =\alpha x_k^{\mathrm{C1}},
                      &                            & \forall k\in\mathcal{K}, \\
    x_k^{\mathrm{D2}} & =\alpha x_k^{\mathrm{D1}},
                      &                            & \forall k\in\mathcal{K}
\end{align}

とする．SOC更新式は

\begin{equation}
    b_{k+1}
    =b_k+
    \left(x_k^{\mathrm{C2}}-x_k^{\mathrm{D1}}\right)\Delta t,
    \quad \forall k\in\mathcal{K}
\end{equation}

とする．蓄電池容量を $C=860$\,kWh，最大充放電出力を $P_{\max}^{\mathrm{C}}=P_{\max}^{\mathrm{D}}=400$\,kW とし，

\begin{align}
    0.05C
        & \leq b_k\leq0.95C,
        &                                    & k=0,\ldots,17520,        \\
    0\leq x_k^{\mathrm{C2}}
        & \leq P_{\max}^{\mathrm{C}}z_k,
        &                                    & \forall k\in\mathcal{K}, \\
    0\leq x_k^{\mathrm{D1}}
        & \leq P_{\max}^{\mathrm{D}}(1-z_k),
        &                                    & \forall k\in\mathcal{K}, \\
    z_k & \in\{0,1\},
        &                                    & \forall k\in\mathcal{K}
\end{align}

を課す．また，年末の蓄電池使い切りによる影響を除くため，初期SOCと終端SOCをともに容量の50\% に固定する．

\begin{equation}
    b_0=b_{17520}=0.5C.
\end{equation}

\subsection{月内最大買電電力と契約電力}

$s_m^{2026}$ は，2026年の月 $m$ における月内最大買電電力なので

\begin{equation}
    s_k^{\mathrm{BY}}\leq s_m^{2026},
    \quad \forall m\in\mathcal{M},\ \forall k\in\mathcal{K}_m
\end{equation}

を満たす．当月を含む直近12か月それぞれの月内最大買電電力を比較し，そのうち最も大きい値を当月の契約電力とする．したがって，月 $m$ の契約電力 $s_m^{\mathrm{CON}}$ は，

\begin{equation}
    \label{eq:contract_rule}
    s_m^{\mathrm{CON}}
    =\max\left(
    \{s_j^{2025}\mid j\in\mathcal{M}_m^{2025}\}
    \cup
    \{s_j^{2026}\mid j\in\mathcal{M}_m^{2026}\}
    \right),
    \quad \forall m\in\mathcal{M}
\end{equation}

と定められる．ただし，式\eqref{eq:contract_rule}は最大値演算を含むため，MILPの線形制約としてそのまま与えることはできない．そこで，$s_m^{\mathrm{CON}}$ を契約電力算定の対象となる全ての月内最大買電電力の上界とみなす．すなわち，2025年に属する月 $j$ の月内最大買電電力 $s_j^{2025}$ と2026年に属する月 $j$ の月内最大買電電力 $s_j^{2026}$ に対して，

\begin{align}
    s_j^{2025} & \leq s_m^{\mathrm{CON}},
               &                          & \forall m\in\mathcal{M},\ \forall j\in\mathcal{M}_m^{2025}, \label{eq:contract_2025} \\
    s_j^{2026} & \leq s_m^{\mathrm{CON}},
               &                          & \forall m\in\mathcal{M},\ \forall j\in\mathcal{M}_m^{2026} \label{eq:contract_2026}
\end{align}

を課す．基本料金単価を $p^{\mathrm{B}}=2{,}405.16$ 円/kW/月，買電電力量単価を $p^{\mathrm{BY}}=21.51$ 円/kWh とし，目的関数を

\begin{equation}
    \label{eq:objective}
    \min\left\{
    \sum_{m\in\mathcal{M}}p^{\mathrm{B}}s_m^{\mathrm{CON}}
    +
    \sum_{k\in\mathcal{K}}p^{\mathrm{BY}}s_k^{\mathrm{BY}}\Delta t
    \right\}
\end{equation}

とする．式\eqref{eq:contract_2025}および式\eqref{eq:contract_2026}だけでは，$s_m^{\mathrm{CON}}$ は各月の月内最大買電電力以上であればよく，最大値より大きな値も許される．しかし，式\eqref{eq:objective}では正の係数 $p^{\mathrm{B}}$ を乗じた $s_m^{\mathrm{CON}}$ を最小化するため，最適解における $s_m^{\mathrm{CON}}$ はこれらの制約を満たす最小値，すなわち当月を含む直近12か月の月内最大買電電力の最大値となる．したがって，式\eqref{eq:contract_2025}，式\eqref{eq:contract_2026}および式\eqref{eq:objective}によって，契約電力算定式\eqref{eq:contract_rule}をMILP上で等価に表している．買電電力を抑制する水準は外部から与えず，契約電力を含む年間電気料金を最小化した結果として決定する．

\section{検証条件}

2025年の月内最大買電電力の大きさと，その値が2026年の契約電力に影響する期間を独立に設定するため，2025年の月内最大買電電力を人工的に作成した．注目する1か月以外の月内最大買電電力は全て100\,kWとし，注目する月のみを140，160，180\,kWのいずれかに設定した．

契約電力は当月を含む直近12か月を参照するため，2025年のどの月に大きな月内最大買電電力を設定するかによって，2026年の契約電力算定に含まれる月数が変化する．具体的な設定を表\ref{tab:remaining_months}に示す．

\begin{table}[H]
    \centering
    \caption{2025年の月内最大買電電力の設定月と契約電力への影響期間}
    \label{tab:remaining_months}
    \begin{tabular}{c p{52mm} p{55mm}}
        \toprule
        影響期間 & 大きな月内最大買電電力を設定する月 & 2026年の契約電力に含まれる月 \\
        \midrule
        1か月  & 2025年2月           & 2026年1月          \\
        6か月  & 2025年7月           & 2026年1--6月       \\
        11か月 & 2025年12月          & 2026年1--11月      \\
        \bottomrule
    \end{tabular}
\end{table}

\begin{table}[H]
    \centering
    \caption{検証条件}
    \label{tab:experiment_conditions}
    \begin{tabular}{ll}
        \toprule
        項目             & 設定値             \\
        \midrule
        2025年の月内最大買電電力 & 140，160，180\,kW \\
        契約電力への影響期間     & 1，6，11か月        \\
        その他の月内最大買電電力   & 100\,kW         \\
        蓄電池容量          & 860\,kWh        \\
        最大充放電出力        & 400\,kW         \\
        充放電効率          & 0.98            \\
        SOC範囲          & 5--95\%         \\
        初期SOC・終端SOC    & 50\%・50\%       \\
        基本料金単価         & 2,405.16円/kW/月  \\
        電力量料金単価        & 21.51円/kWh      \\
        \bottomrule
    \end{tabular}
\end{table}

2025年の月内最大買電電力3水準と契約電力への影響期間3水準の組合せによる9ケースを計算した．

\section{計算結果}

9ケースのうち，月内最大買電電力と影響期間が小・中・大となる3条件を図\ref{fig:representative_conditions}に示す．横軸は2025年の月，縦軸は各月の月内最大買電電力である．各段はそれぞれ別の計算条件を表し，注目月以外の月内最大買電電力は100\,kWに設定した．

\begin{figure}[H]
    \centering
    \begin{tikzpicture}[x=0.66cm,y=0.025cm]
        \foreach \offset/\peakmonth/\peakvalue/\duration in {5.4/2/140/1,2.7/7/160/6,0/12/180/11}{
                \begin{scope}[yshift=\offset cm]
                    \draw[->] (0,0) -- (12.4,0);
                    \draw[->] (0,0) -- (0,90);
                    \foreach \tick/\label in {0/100,40/140,80/180}{
                            \draw (-0.12,\tick) -- (0.12,\tick);
                            \node[anchor=east] at (-0.22,\tick) {\scriptsize \label};
                        }
                    \draw[densely dashed,black!45] (0.5,0) -- (11.5,0);
                    \foreach \m in {1,...,12}{
                            \draw ({\m-0.5},-3) -- ({\m-0.5},3);
                            \node[anchor=north] at ({\m-0.5},-5) {\tiny \m};
                            \ifnum\m=\peakmonth\relax
                            \else
                                \fill ({\m-0.5},0) circle (1.2pt);
                            \fi
                        }
                    \draw[thick] ({\peakmonth-0.5},0) -- ({\peakmonth-0.5},{\peakvalue-100});
                    \fill ({\peakmonth-0.5},{\peakvalue-100}) circle (2pt);
                    \node[anchor=south] at ({\peakmonth-0.5},{\peakvalue-100}) {\small \peakvalue\,kW};
                    \node[anchor=west] at (12.65,40) {\small 影響期間：\duration か月};
                \end{scope}
            }
        \node[rotate=90] at (-1.4,190) {月内最大買電電力 [kW]};
        \node at (6.2,-35) {2025年の月};
    \end{tikzpicture}
    \caption{代表3条件における2025年の月内最大買電電力}
    \label{fig:representative_conditions}
\end{figure}

9ケースの年間基本料金，年間電力量料金および年間電気料金に加え，年間最大契約電力を表\ref{tab:cost_results}に示す．

\begin{table}[H]
    \centering
    \small
    \caption{2025年の月内最大買電電力と影響期間による年間電気料金}
    \label{tab:cost_results}
    \begin{tabular}{rrrrrr}
        \toprule
        \shortstack{2025年の月内                                                          \\最大買電電力}
                             & 影響期間
                             & \shortstack{年間最大                                       \\契約電力}
                             & \shortstack{年間                                         \\基本料金}
                             & \shortstack{年間                                         \\電力量料金}
                             & \shortstack{年間                                         \\電気料金} \\
        {[kW]}               & {[月]}            & {[kW]} & {[万円]} & {[万円]}  & {[万円]}  \\
        \midrule
        \multirow{3}{*}{140} & 1                & 166.8  & 396.8  & 1,136.5 & 1,533.3 \\
                             & 6                & 166.8  & 439.6  & 1,136.5 & 1,576.1 \\
                             & 11               & 166.8  & 439.6  & 1,136.5 & 1,576.1 \\
        \midrule
        \multirow{3}{*}{160} & 1                & 166.8  & 401.6  & 1,136.5 & 1,538.1 \\
                             & 6                & 166.8  & 468.5  & 1,136.4 & 1,604.9 \\
                             & 11               & 166.8  & 470.0  & 1,136.3 & 1,606.3 \\
        \midrule
        \multirow{3}{*}{180} & 1                & 180.0  & 406.4  & 1,136.5 & 1,542.9 \\
                             & 6                & 180.0  & 497.3  & 1,136.4 & 1,633.8 \\
                             & 11               & 180.0  & 516.3  & 1,136.2 & 1,652.6 \\
        \bottomrule
    \end{tabular}
\end{table}

9ケースについて，2026年各月の月内最大買電電力，契約電力および月間電気料金を図\ref{fig:monthly_peak140_remain1}--\ref{fig:monthly_peak180_remain11}に示す．月間電気料金は，各月の基本料金と電力量料金の和である．

\newcommand{\monthlycasefigure}[2]{%
    \begin{figure}[H]
        \centering
        \begin{tikzpicture}
            \begin{axis}[
                width=0.90\textwidth,
                height=6.2cm,
                xmin=1, xmax=12,
                ymin=80, ymax=190,
                xtick={1,...,12},
                ytick={80,100,120,140,160,180},
                xlabel={2026年の月},
                ylabel={電力 [kW]},
                axis y line*=left,
                axis x line*=bottom,
                grid=major,
                grid style={black!15},
                tick label style={font=\small},
                label style={font=\small},
                legend style={
                        at={(0.5,-0.25)},
                        anchor=north,
                        legend columns=3,
                        draw=none,
                        font=\small
                    }
                ]
                \addplot[black,thick,mark=*]
                table[
                        x expr=\coordindex+1,
                        y=contract_power_kW,
                        col sep=comma
                    ] {../results/demand_target_sensitivity/cases/peak#1_remain#2_targetnone/monthly.csv};
                \addlegendentry{契約電力}
                \addplot[black,thick,dashed,mark=square*]
                table[
                        x expr=\coordindex+1,
                        y=monthly_peak_kW,
                        col sep=comma
                    ] {../results/demand_target_sensitivity/cases/peak#1_remain#2_targetnone/monthly.csv};
                \addlegendentry{月内最大買電電力}
                \addlegendimage{black,thick,densely dotted,mark=triangle*}
                \addlegendentry{月間電気料金}
            \end{axis}
            \begin{axis}[
                width=0.90\textwidth,
                height=6.2cm,
                xmin=1, xmax=12,
                ymin=0, ymax=300,
                xtick=\empty,
                ytick={0,50,100,150,200,250,300},
                ylabel={月間電気料金 [万円]},
                axis y line*=right,
                axis x line=none,
                tick label style={font=\small},
                label style={font=\small}
                ]
                \addplot[black,thick,densely dotted,mark=triangle*]
                table[
                        x expr=\coordindex+1,
                        y expr=\thisrow{total_cost_yen}/10000,
                        col sep=comma
                    ] {../results/demand_target_sensitivity/cases/peak#1_remain#2_targetnone/monthly.csv};
            \end{axis}
        \end{tikzpicture}
        \caption{2025年の月内最大買電電力#1\,kW・影響期間#2か月の条件における2026年の月別推移}
        \label{fig:monthly_peak#1_remain#2}
    \end{figure}
}

\monthlycasefigure{140}{1}
\monthlycasefigure{140}{6}
\monthlycasefigure{140}{11}
\monthlycasefigure{160}{1}
\monthlycasefigure{160}{6}
\monthlycasefigure{160}{11}
\monthlycasefigure{180}{1}
\monthlycasefigure{180}{6}
\monthlycasefigure{180}{11}

2025年の月内最大買電電力が大きく，その値が契約電力算定に含まれる期間が長いほど，年間基本料金はおおむね増加した．

月内最大買電電力を140\,kWとした場合は，影響期間6か月と11か月の年間電気料金が同じであった．影響期間6か月の条件では，2026年7月以降に140\,kWを上回る月内最大買電電力が生じる．このため，140\,kWが契約電力算定にさらに5か月残っていても契約電力を決める値とはならず，年間電気料金に追加の影響を与えなかった．

計算にはPySCIPOpt 5.5.0，SCIP 9.2.1を用いた．9ケースのソルバ計算時間の合計は617.6秒，1ケースの最大計算時間は74.3秒であった．

\section{まとめ}

本検証により，2025年の月内最大買電電力は，その大きさだけでなく，契約電力算定に含まれる期間によって2026年の年間電気料金に与える影響が変化することを確認した．月内最大買電電力が140，160，180\,kW，影響期間が1，6，11か月の9ケースでは，年間電気料金に最大約119.3万円の差が生じ，その大部分は年間基本料金の差であった．

\end{document}
